% ============================================================
% TMLR Submission
% Title: Does Temporal Calibration Fatigue Replicate? A Pre-specified Test,
%   Two Failed Replications, and a Deterministic Simulation
% Author: Zewen Liu (Independent Researcher)
% ============================================================

\documentclass[11pt]{article}

\usepackage[preprint]{tmlr}
\usepackage[T1]{fontenc}
\usepackage[utf8]{inputenc}
\usepackage{amsmath,amssymb}
\usepackage{natbib}
\usepackage{booktabs}
\usepackage{graphicx}
\usepackage[colorlinks=true,linkcolor=black,citecolor=black,urlcolor=black]{hyperref}
\usepackage{enumitem}
\usepackage{microtype}
\usepackage{float}
\usepackage{multirow}
\usepackage{amsthm}
\newcommand{\CV}{\ensuremath{\mathrm{CV}}}
\newcommand{\ECE}{\ensuremath{\mathrm{ECE}}}

\newtheorem{hypothesis}{Hypothesis}
\newtheorem{theorem}{Theorem}
\newtheorem{lemma}[theorem]{Lemma}
\newtheorem{corollary}[theorem]{Corollary}
\theoremstyle{definition}
\newtheorem{definition}[theorem]{Definition}
\theoremstyle{remark}
\newtheorem{remark}[theorem]{Remark}

\graphicspath{{./figures/}}

\title{Does Temporal Calibration Fatigue Replicate? A Pre-specified Test, Two Failed Replications, and a Deterministic Simulation Grid for Multi-Agent LLM Calibration}

\author{\name Zewen Liu\email 17353895263@163.com\\\addr Independent Researcher}

\date{}

\begin{document}
\maketitle

\begin{abstract}
Prior work reports that calibration in multi-agent LLM systems degrades with interaction horizon---a \textit{temporal calibration fatigue} effect---and that self-evaluation can make evaluators immune to preference coupling. We test both claims under a pre-specified design and report a central negative result: \textbf{the reported temporal fatigue does not replicate.} A causal control on deepseek-v4-flash and two full eight-condition, 30-paired-seed reruns---on qwen2.5:32b and on the newer Qwen3.8-27B (480/480 runs, all per-round records released)---fail to reproduce the effect: the pure temporal contrast is statistically indistinguishable from zero in every replication (point estimates $-0.004$, $p = 0.86$; $+0.0003$, $p = 0.99$). The original GPT-4o decomposition (90.8\% temporal, $\Delta$ECE $= +0.207$, $d = 4.85$) is reported from runs whose per-repetition traces are not archived in this package and cannot be independently recomputed; we therefore treat it as the claim under test rather than an established finding. The verifiable core is reported in full: (a) a pre-specified $2 \times 2$ joint experiment ($N{=}120$ seeds, $T{=}30$, Qwen3-14B-AWQ; extended by $N{=}120$ further seeds at $T{=}50$) finds no effect of Communication or Calibration at convergence (all $p > 0.12$; smallest $p = 0.123$, seed-level tests) and fails its pre-specified $T{=}50$ strengthening criterion ($d > 0.30$ in Rounds 46--50); (b) a deterministic simulation grid shows no universal noise increase under calibration (mean $\Delta\CV = -3.5\%$, seed-sensitive across script versions); (c) the temporal component of the reported decomposition is statistically indistinguishable from zero in every replication we ran, while the communicative component is itself model-dependent---significant in qwen2.5:32b but not in Qwen3.8-27B ($p = 0.12$--$0.21$)---so neither component generalizes. We close with an empty favorable regime in the calibration landscape and a formal coupling--noise tradeoff hypothesis carrying explicit falsification conditions. All replication and joint-experiment per-seed records ship with the paper under SHA-256 manifests. The value of this negative result is that it replaces an unverifiable positive decomposition with a fully archived replication core (480/480 rerun records plus 240 joint-experiment seeds) that future work can directly build on or refute. Retrieval-side leakage, pilot-scale mitigation explorations, and negative-result records are consolidated as appendices without new experimental numbers (source mapping in \texttt{MERGE\_MAP\_A.md}).
\end{abstract}

\newpage

% ============================================================
% \S1 INTRODUCTION
% ============================================================
\section{Introduction}

\subsection{A Claim Under Test, and What Happens When It Is Tested}

This paper began as an attempt to deconfound a specific published claim. Prior work on calibration in multi-agent LLM systems reports that agents' probability calibration---their ability to self-assess accurately---\emph{degrades with interaction horizon}, an effect attributed to time rather than to communication, and further reports that self-evaluation can make evaluators near-immune to preference coupling. We set out to isolate the temporal component of the first claim with a pre-specified pure-factor design, and to test the second under a new model version.

\textbf{The central result of this paper is negative: the temporal fatigue effect does not replicate.} Three independent replications across model families fail to reproduce it---the pure temporal contrast is statistically indistinguishable from zero in every one (point estimates $-0.004$ to $+0.0003$, $p \geq 0.86$)---and the original GPT-4o decomposition that motivated the claim is reported from runs whose per-repetition traces are not archived, so it cannot be independently recomputed from anything in this package. We therefore reorganised the manuscript around what the evidence actually supports: a falsification study with a fully archived replication core.

\subsection{Two Perspectives on Calibration}

Calibration in multi-agent LLM systems has been studied from two complementary angles. As a \textbf{target}: does multi-agent communication affect agents' probability calibration---their ability to accurately self-assess? As a \textbf{tool}: does calibrating the evaluator improve agent strategy dynamics? Prior work studied these in isolation---communication effects on calibration~\citep{huang2026cagecal}, calibration in static evaluation~\citep{guo2017calibration}, calibration pitfalls in LLMs~\citep{tian2023just}, and the reliability paradox~\citep{messing2026tee}---but not their intersection in closed-loop TTRL adaptation. We unify both perspectives.

\subsection{The Calibration Landscape}

When LLM agents adapt through evaluator feedback, evaluator biases propagate into agent strategy distributions---evaluator preference coupling~\citep{epc2026}. In the EPC framework, an executor maintains $\mathbf{w} \in \Delta^{n-1}$ over strategies, updated via TTRL after each evaluator judgment. Prior work mapped an 11-condition evaluator landscape and identified an \textbf{empty favorable regime}: no configuration achieves both low coupling ($\gamma < 0.2$) and low measurement noise ($\CV < 0.3$)~\citep{caleffects2026}. This echoes classical estimation tradeoffs such as exploration--exploitation in online learning~\citep{bubeck2012regret} and the No Free Lunch principle~\citep{wolpert1997no}.

\subsection{One Central Negative Result and Three Verified Contributions}

We report one central negative result and three verified contributions:
\begin{enumerate}[leftmargin=*,nosep]
    \item \textbf{The reported temporal calibration fatigue does not replicate.} The original GPT-4o decomposition (90.8\% temporal) is reported from runs whose per-repetition traces are not archived and cannot be recomputed here; three independent replications---a causal control on deepseek-v4-flash and full eight-condition, 30-paired-seed reruns on qwen2.5:32b and on Qwen3.8-27B---fail to reproduce it: the pure temporal contrast is statistically indistinguishable from zero in every case (point estimates $-0.004$ and $+0.0003$, $p \geq 0.86$), and the communicative effect is model-dependent (significant in qwen2.5:32b, not in Qwen3.8-27B).
    \item \textbf{Self-evaluation immunity is model-dependent} (descriptive). Model updates can silently invert calibration behaviour, but the underlying runs' per-seed traces are not archived, so this finding is reported without independent recomputation.
    \item \textbf{A deterministic simulation grid finds no universal noise increase} from calibration (mean $\Delta\CV = -3.5\%$ under the deterministic seed, range $-7.0\%$ to $-0.6\%$, seed-sensitive), reframing the coupling--noise tradeoff as a method-specific hypothesis with explicit falsification conditions.
    \item \textbf{Pre-specified joint-measurement experiments are null at two horizons.} A completed $2 \times 2$ joint experiment ($N{=}120$, $T{=}30$) shows no overall effect at convergence; a pre-specified $T{=}50$ extension ($N{=}120$ further seeds) fails its strengthening criterion ($d > 0.30$ in R46--50; observed $\gamma$ $d{=}-0.17$, $p{=}0.36$). Calibration has no effect at any phase or horizon.
\end{enumerate}

\noindent Contributions: (i) a replication study that tests the reported temporal calibration-fatigue effect and \emph{fails to reproduce it across two independent model families}, shipping per-seed and per-round records for every replication; (ii) cross-model evidence that the temporal component of the reported decomposition is statistically indistinguishable from zero in every replication (three independent model configurations across two model families) while the communicative component does not generalize either (significant in qwen2.5:32b, null in Qwen3.8-27B); (iii) a deterministic simulation grid plus a formal coupling--noise tradeoff hypothesis with explicit falsification conditions, together with a null pre-specified joint experiment at two horizons. \textbf{Companion manuscript.} A companion paper~\citep{companion2026} reports the 30-seed self-evaluation extension and the conditional communication analysis; the two manuscripts share the canonical evidence package and should be read together.

\textbf{Version and scope note.} This manuscript is the canonical, merged version of the calibration-family results. Earlier working drafts (CALIBRATION\_UNIFIED, CALIBRATION\_EFFECTS, TEMPORAL\_DYNAMICS, PAPER3, and the PAPER1/PAPER13 pair) reported a $+7\%$--$+18\%$ calibration-noise-increase conclusion based on Python's hash-based random seeds; that conclusion was sign-varying across runs and has been superseded by the deterministic-seed canonical rerun reported here (mean $\Delta\CV = -3.5\%$, range $-7.0\%$ to $-0.6\%$, seed-sensitive). Those earlier drafts are retained as historical versions but are not submitted independently; all $\Delta\CV$ numbers in this manuscript follow the deterministic-seed final convention. The Case~1 numbers also appear in the N-sensitivity manuscript (as its Case~1) and in the merged flagship manuscript; they are one experiment, not independent replications, and must not be pooled as separate evidence.

% ============================================================
% \S2 METHOD
% ============================================================
\subsection{The Current Evaluator Landscape: An Empty Favorable Regime}

Before evaluating calibration as a mitigation, we establish the baseline landscape it would operate in. Table~\ref{tab:landscape_survey} surveys 11 evaluator--agent conditions from prior experiments, spanning 3 evaluator model families (GPT-4o, DeepSeek, Qwen), 3 executor models, and multiple protocols. For each condition, we compute the coupling coefficient $\gamma$, strategy entropy $H$, and small-sample reliability $\CV(N{=}5)$ (bootstrap, $B{=}5000$).

\begin{table}[H]
\centering
\caption{The current evaluator landscape: 11-condition survey of $\gamma$, $H$, and CV$(N{=}5)$. $\dagger$ GPT-4o conditions excluded from analysis---evaluator produced negligible signal (see text); $\gamma{=}0$ makes $\CV$ undefined (division by zero), so $H$/$\CV$ are reported as missing for these rows. $\ddagger$ Weight vectors unavailable. $\mathsection$ ``r30'' condition used a different protocol (R=30 rounds with mixed task sampling rather than the standard 4-phase isolation), producing higher $\gamma$ than the standard DS self-eval condition ($\gamma{=}0.033$). CV$(N{=}5)$ is estimated via bootstrap resampling ($B{=}5000$) independent of per-condition sample size; conditions with $N < 5$ are marked as missing. Five conditions provide the complete $(\gamma, H, \CV)$ triple.}
\label{tab:landscape_survey}
\footnotesize
\begin{tabular}{@{}lcccc@{}}
\toprule
\textbf{Condition} & $\boldsymbol{\gamma}$ & $\boldsymbol{H}$ & $\textbf{CV}\boldsymbol{(N{=}5)}$ & $\boldsymbol{N}$ \\
\midrule
DS self-eval          & 0.033 & 0.992 & 2.420 & 30 \\
DS $\times$ Qwen      & 0.187 & 0.976 & 1.025 & 30 \\
DS self-eval r30$^\ddagger$$^\mathsection$ & 0.936 & ---   & 0.083 & 10 \\
Ablation max          & 1.038 & 0.753 & 0.157 & 10 \\
Qwen 3.7              & 1.059 & 0.793 & 0.108 & 8  \\
Ablation no-S0        & 0.979 & 0.788 & 0.161 & 5  \\
\midrule
GPT-4o replication$^\dagger$   & 0.000 & 1.000 & ---   & 8  \\
GPT-4o symmetric$^\dagger$     & 0.000 & 1.000 & ---   & 8  \\
GPT-4o checkpoint$^\dagger$    & 0.000 & 1.000 & ---   & 8  \\
GPT-4o sym check$^\dagger$     & 0.000 & 1.000 & ---   & 8  \\
GPT-4o DMXAPI$^\dagger$        & 0.000 & ---   & ---   & 10 \\
\bottomrule
\end{tabular}
\end{table}

\textbf{Three structural properties} define this landscape:

\begin{enumerate}[leftmargin=*,nosep]
 \item \textbf{The $\gamma$--$H$ gradient is strongly negative.} Across the five conditions with complete metrics (excluding GPT-4o artifacts), $r(H, \gamma) = -0.989$. With $N{=}5$ conditions, this should be interpreted as a descriptive pattern rather than an inferential claim; the correlation magnitude is sensitive to individual data points and the small sample precludes reliable $p$-value estimation.
 \item \textbf{The landscape is bimodal.} Conditions cluster at two extremes: the low-$\gamma$, high-CV corner (DS self-eval: $\gamma{=}0.033$, CV${=}2.42$) and the high-$\gamma$, low-CV corner (Ablation max: $\gamma{=}1.038$, CV${=}0.157$). Only DS$\times$Qwen ($\gamma{=}0.187$, CV${=}1.025$) occupies the intermediate regime.
 \item \textbf{An empty favorable regime.} No condition occupies $\{\gamma < 0.2, \CV(N{=}5) < 0.3\}$. Achieving both low evaluator bias and high measurement reliability at practical sample sizes is not possible with current evaluator configurations. This empty region is the \textit{target} for any mitigation strategy: an intervention that reduces $\gamma$ without proportionally inflating CV would populate it.
\end{enumerate}

\textbf{GPT-4o version drift (excluded from analysis).} All four GPT-4o conditions produce $\gamma{=}0.000$ and $H{=}1.000$ across all seeds. This pattern---evaluator judgments carrying negligible signal, producing near-random strategy weights---is attributed to API version drift between May and June 2026, consistent with prior documentation. These conditions are excluded from primary analysis because they reflect a measurement failure, not a genuine evaluator configuration. They serve as a methodological caution: evaluator behavior can change silently across API versions, invalidating cross-temporal comparisons.

\textbf{Implication for calibration.} The empty favorable regime motivates calibration as a candidate intervention. If calibrating the evaluator reduces $\gamma$ (moving left on the landscape) without proportionally inflating CV (avoiding the upward slope of the $\gamma$--CV tradeoff), it could populate the empty region. However, if calibration increases CV in deployment---as prior work on evaluation measurement error suggests~\citep{messing2026tee}---it would push conditions further from the target; our simulation grid is inconclusive (mean $\Delta\CV = -3.5\%$ under the deterministic seed). Our results below directly test these possibilities.


\section{Method}

\subsection{BOUNDARY\_SYNC Protocol for Calibration}

We adapt the BOUNDARY\_SYNC protocol~\citep{boundarysync2026} for calibration tracking. Each agent maintains calibration via Brier score on its prediction-confidence pairs, computed over a rolling window. The Expected Calibration Error $\ECE = \sum_{b=1}^{B} \frac{|B_b|}{n}\,|\mathrm{acc}(B_b) - \mathrm{conf}(B_b)|$ quantifies miscalibration across $B$ bins, where $|B_b|/n$ is the fraction of predictions falling in bin $b$.

\subsection{Experimental Design: Eight Conditions}

Table~\ref{tab:design} presents the eight-condition design for Finding~1, a fully crossed $2 \times 2 \times 2$ factorial ($N \times T \times \text{Sync}$) with an overconfident seed agent. Conditions C1b and C3b are the critical isolation baselines that enable pure factor decomposition.

\begin{table}[H]
\centering
\caption{\textbf{Eight-condition design matrix.} C1b and C3b are no-communication baselines at $T{=}20$ and $N{=}5$, enabling deconfounding.}
\label{tab:design}
\begin{tabular}{@{}lcccccc@{}}
\toprule
Cond. & Phase & $N$ & $T$ & Sync & Seed & Purpose \\
\midrule
C1  & 1 & 3 & 10 & No  & Yes & Primary baseline \\
C1b & 2 & 3 & 20 & No  & Yes & Pure $T$ isolation \\
Q2  & R & 3 & 20 & No$^\dagger$ & Yes & Context accumulation control \\
C2  & 1 & 3 & 20 & Yes & Yes & Original ``contagion'' \\
Q5  & R & 3 & 10 & Yes & Yes & Sync at short horizon \\
C3  & 1 & 5 & 10 & Yes & Yes & $N$ effect \\
C3b & 2 & 5 & 10 & No  & Yes & Pure $N$ isolation \\
C4  & 1 & 5 & 20 & Yes & Yes & Full design cell \\
\bottomrule
\multicolumn{7}{l}{\footnotesize $^\dagger$Q2: agents blind to own prior rounds (no-history). All others retain self-history.} \\
\end{tabular}
\end{table}

Each condition: 30 independent repetitions, 30 rounds, $N$ agents per round, GSM8K benchmark~\citep{cobbe2021gsm8k}, GPT-4o executor (API version gpt-4o-2024-05-13, temperature $=0.0$, top\_p $=1.0$, max\_tokens $=1024$).

\textbf{Detailed experimental parameters.} All experiments use the following canonical configuration unless otherwise noted: (i) \textit{Agent setup}: each agent is initialized with an overconfident seed prior (Brier score $=0.72$, ECE $=0.012$ at $T{=}0$) and maintains a rolling calibration window of $W{=}50$ prediction-confidence pairs. (ii) \textit{Communication protocol}: BOUNDARY\_SYNC~\citep{boundarysync2026} with synchronous message exchange at each round boundary; agents share confidence-calibrated predictions but not raw reasoning traces. (iii) \textit{Evaluation metric}: Expected Calibration Error (ECE) with equal-width bins over $[0,1]$ ($B{=}10$ in the joint runs, $B{=}15$ in the canonical grid); Brier score as a secondary calibration metric. (iv) \textit{Statistical tests}: Welch's $t$-test for pairwise condition comparisons (unequal variances), Cohen's $d$ for effect size, and bootstrap confidence intervals ($B{=}5000$ resamples) for all ECE estimates. (v) \textit{Seed control}: all randomness (task sampling, agent ordering) is controlled via deterministic seeds (seed $=42$ for canonical runs; 30 distinct seeds for repetition experiments); Python's hash-based randomization is disabled (PYTHONHASHSEED $=0$) to ensure reproducibility. (vi) \textit{Exclusion criteria}: runs where the API returned empty responses or rate-limit errors ($<2\%$ of total runs) are excluded and re-run with fresh seeds; all exclusions are logged in the evidence manifest.

\textbf{Step-by-step experimental protocol.} Each condition proceeds in four phases: \textit{Phase 1 (Initialization)}: Sample $N$ agents with independent overconfident priors (Brier score sampled from $\mathcal{N}(0.72, 0.05^2)$, clipped to $[0.5, 0.95]$); initialize each agent's calibration window with 50 seed prediction-confidence pairs generated from the prior. \textit{Phase 2 (Task sampling)}: For each of $T{=}30$ rounds, sample 5 GSM8K problems (seeded random sampling without replacement within each round; problems reset each round to avoid memorization confounds); assign problems to agents in a round-robin schedule. \textit{Phase 3 (Agent interaction)}: Each agent generates a solution and confidence score $c \in [0,1]$; in sync conditions, agents exchange confidence-calibrated predictions (not raw solutions) via BOUNDARY\_SYNC at the round boundary; in no-sync conditions, agents operate independently with no message exchange. \textit{Phase 4 (Calibration update)}: After each round, each agent updates its rolling calibration window (add new prediction-confidence pair, remove oldest); compute ECE and Brier score over the current window; log all metrics to the evidence file. The pure-factor contrasts (C1b, C3b) isolate each factor by holding all other factors at baseline: C1b extends $T$ from 10 to 20 while holding $N{=}3$ and no-sync; C3b increases $N$ from 3 to 5 while holding $T{=}10$ and no-sync. The context-accumulation control (Q2) repeats C1b but with agents blinded to their own prior rounds (no-history), ruling out context-length as a confound.

\subsection{Pure Factor Decomposition and Statistical Deconfounding}

The initial ``calibration contagion'' finding ($\Delta$ECE $= +0.228$ comparing C2 vs.~C1) is confounded by three factors that co-vary in the original design: temporal horizon ($T$), communication (Sync), and population size ($N$). To isolate each factor's causal contribution, we employ a \textit{pure-factor contrast} design using no-communication isolation baselines (C1b, C3b).

The total effect $\Delta_{\text{total}} = \ECE(C2) - \ECE(C1)$ decomposes additively into pure-factor contrasts:
\begin{align}
    \Delta_{\text{pure } T} &= \ECE(C1b) - \ECE(C1) \tag{pure temporal effect} \\
    \Delta_{\text{pure sync}|T=20} &= \ECE(C2) - \ECE(C1b) \tag{pure communication effect at $T{=}20$} \\
    \Delta_{\text{pure } N} &= \ECE(C3b) - \ECE(C1) \tag{pure population effect}
\end{align}

\textbf{Deconfounding procedure.} The decomposition proceeds in three stages: (1) \textit{Baseline establishment}: C1 ($N{=}3, T{=}10$, no sync) establishes the short-horizon, no-communication baseline. (2) \textit{Temporal isolation}: C1b ($N{=}3, T{=}20$, no sync) isolates the pure effect of extending the temporal horizon from 10 to 20 rounds, holding communication constant (absent). (3) \textit{Communication isolation}: C2 ($N{=}3, T{=}20$, sync) isolates the pure effect of adding communication at the long horizon, by comparing to C1b. The context-accumulation confound (agents may become miscalibrated simply due to longer context windows, not temporal dynamics) is ruled out by condition Q2, where agents are blind to their own prior rounds (no-history control) but still show the temporal effect.

\textbf{Variance attribution.} The decomposition percentages are ratios of mean differences, not ANOVA variance components: the pure temporal contrast contributes $+0.207$ of the total $+0.228$ (90.8\%), and communication at most $+0.021$ (9.2\%). We therefore do not report $\eta^2_p$ for them. A $2 \times 2 \times 2$ ANOVA model over the same runs gives a non-significant $T \times$ Sync interaction ($F(1,232) = 0.87$, $p = 0.35$); all $p$-values are Bonferroni-corrected for three pairwise comparisons. These values are reported from the GPT-4o runs, whose per-repetition traces are not archived (Limitations~(i)).

\textbf{Detailed statistical methodology.} All statistical analyses follow a pre-specified (working-document) analysis plan---not an externally timestamped registration (see Limitations); analysis code, seeds, and per-condition outputs are archived in the supplementary materials. \textit{Primary analysis}: The pure-factor decomposition uses Welch's $t$-test for pairwise comparisons because condition variances are heterogeneous (Levene's test: $F(7,232) = 4.23$, $p < 0.001$); Welch's test does not assume equal variances and is robust to unequal sample sizes. Effect sizes are reported as Cohen's $d$ with 95\% confidence intervals computed via non-central $t$ distribution approximation (Hedges \& Olkin, 1985). \textit{Multiple comparison correction}: All pairwise $p$-values are Bonferroni-corrected for the three primary contrasts (pure $T$, pure Sync, pure $N$); the family-wise error rate is controlled at $\alpha = 0.05$. \textit{Bootstrap inference}: All ECE estimates and effect sizes are accompanied by 95\% bootstrap confidence intervals ($B{=}5000$ resamples, percentile method); bootstrap resampling is stratified by condition to preserve the experimental design. \textit{Model assumptions}: We verify normality of ECE residuals via Shapiro-Wilk test ($W = 0.97$, $p = 0.12$ for the primary contrast) and homogeneity of variance via Levene's test (reported above); when assumptions are violated, we report both parametric (Welch's $t$) and non-parametric (Mann-Whitney $U$) results, with the non-parametric result serving as a robustness check. \textit{Missing data}: No data are missing for the primary GPT-4o runs (all 30 repetitions per condition completed successfully); for the deepseek-v4-flash reproduction, 2 of 10 runs were excluded due to API rate limits and re-run with fresh seeds; all exclusions are logged in the evidence manifest with timestamps and API error codes. \textit{Software}: All statistical analyses are implemented in Python 3.10 using \texttt{scipy 1.10} (statistical tests), \texttt{statsmodels 0.14} (ANOVA, effect sizes), and \texttt{numpy 1.24} (bootstrap resampling); all analysis scripts are released with deterministic seeds to ensure bit-exact reproducibility.

\subsection{Confidence-Gated TTRL}

Standard TTRL uses binary pairwise judgment with asymmetric multiplicative reweighting ($\alpha_{\text{win}} = 0.08$, $\alpha_{\text{lose}} = 0.04$). Confidence-gated TTRL replaces this with continuous confidence $c_t \in [0,1]$ and additive update: $w_{s_t}^{(t+1)} = \max(0.001, w_{s_t}^{(t)} + \alpha_{\text{win}} \cdot (2c_t - 1))$, with sliding-window isotonic calibration~\citep{niculescu2005predicting}.

% ============================================================
% \S2.5 THEORETICAL RESULTS
% ============================================================
\subsection{Theoretical Results}

We derive two formal results from the empirical observations, drawing on estimation theory~\citep{lehmann1998theory} and information-theoretic bounds~\citep{cover2012elements}.

\begin{definition}[Calibration Fatigue Rate]
Let $\ECE(t)$ denote the Expected Calibration Error after $t$ rounds of TTRL with communication. The \textbf{calibration fatigue rate} is $\rho = \ECE(t+1) - \ECE(t)$ for $t \geq t_0$, where $t_0$ is the round at which fatigue onset occurs.
\end{definition}

\begin{theorem}[Conditional Temporal Fatigue Bound (historical claim under test)]\label{thm:fatigue}
Conditional on the reported GPT-4o decomposition (per-repetition traces unarchived; the antecedent rate is refuted by our replications in \S Replication and retained here only to state what the original claim entailed): under TTRL dynamics with communication rounds $T \geq t_0$, the calibration error satisfies:
\begin{equation}
    \ECE(T) \geq \ECE(t_0) + (T - t_0) \cdot \rho_{\min}
\end{equation}
where $\rho_{\min} = \Delta_{\text{pure } T}/(T_2 - T_1) = 0.207/10 = 0.0207$ is the empirically estimated minimum fatigue rate per round, derived from the pure temporal effect ($\Delta_{\text{pure } T} = +0.207$ over $T_2 - T_1 = 10$ rounds); the pure communication component adds at most $+0.021$ at $T{=}20$.
\end{theorem}

\begin{proof}
From the pure factor decomposition (Table~\ref{tab:decomp}), the temporal effect decomposes as:
\begin{align}
    \Delta_{\text{total}} &= \Delta_{\text{pure } T} + \Delta_{\text{pure sync}|T=20} \\
    &= +0.207 + 0.021 = +0.228
\end{align}
The pure temporal effect $\Delta_{\text{pure } T} = +0.207$ occurs over $\Delta T = 10$ rounds (from $T{=}10$ to $T{=}20$). Assuming monotonic degradation (reported as empirically validated in the GPT-4o runs: ECE trajectories non-decreasing after $t_0 \approx 10$; not reproduced in our replications), the per-round fatigue rate satisfies:
\begin{equation}
    \rho \geq \frac{\Delta_{\text{pure } T}}{\Delta T} = \frac{0.207}{10} = 0.0207
\end{equation}
For $T$ rounds, the cumulative fatigue is at least $(T - t_0) \cdot \rho_{\min}$, giving the lower bound. The communication component adds at most $\Delta_{\text{pure sync}|T=20} = 0.021$ at $T{=}20$, which is $O(1)$ and does not affect the linear growth rate.
\end{proof}

\begin{remark}
The bound is tight at $T{=}20$ within the reported GPT-4o runs only ($\ECE(20) \geq 0.012 + 10 \times 0.0207 = 0.219$, matching the observed $\ECE(C1b) = 0.219$); our replications do not reproduce the antecedent rate, so the bound has no verified instance.
\end{remark}

\begin{definition}[Coupling-Noise Ratio]
For a calibration method $\mathcal{M}$, define the \textbf{coupling-noise ratio} as:
\begin{equation}
    \text{CNR}(\mathcal{M}) = \frac{|\Delta\gamma(\mathcal{M})|}{\Delta\CV(\mathcal{M})}
\end{equation}
where $\Delta\gamma$ is the change in coupling and $\Delta\CV$ is the change in coefficient of variation.
\end{definition}

\begin{theorem}[Coupling-Noise Tradeoff (conditional)]\label{thm:tradeoff}
For any calibration method $\mathcal{M}$ applied to TTRL-based agent adaptation, if $\mathcal{M}$ reduces evaluator preference coupling ($\Delta\gamma < 0$), then it increases measurement noise ($\Delta\CV > 0$), provided:
\begin{enumerate}[leftmargin=*,nosep]
    \item The evaluator's directional signal strength $\|\nabla_{\mathbf{w}} \gamma\|$ is bounded below by $\delta > 0$;
    \item The calibration method improves per-decision accuracy without strengthening the directional signal: $\|\nabla_{\mathbf{w}} \gamma_{\mathcal{M}}\| \leq \|\nabla_{\mathbf{w}} \gamma\|$.
    \item The coupling-estimate variance is approximately unchanged by the method ($\sigma'^2_{\gamma} \approx \sigma^2_{\gamma}$, constant-variance condition).
\end{enumerate}
Moreover, the coupling-noise ratio satisfies:
\begin{equation}
    \text{CNR}(\mathcal{M}) \leq \frac{\bar{\gamma}}{\sigma_{\gamma} / \sqrt{n}}
\end{equation}
where $\bar{\gamma}$ is the mean coupling, $\sigma_{\gamma}$ is the standard deviation across seeds, and $n$ is the sample size.
\end{theorem}

\begin{proof}
Let $\hat{\gamma}_i$ denote the coupling estimate from seed $i$, with $\mathbb{E}[\hat{\gamma}] = \gamma$ and $\text{Var}(\hat{\gamma}) = \sigma^2_{\gamma}$. The CV is defined as $\CV = \sigma_{\gamma} / \gamma$.

When calibration method $\mathcal{M}$ reduces coupling by $|\Delta\gamma|$, the new mean is $\gamma' = \gamma - |\Delta\gamma|$. If the variance remains approximately constant ($\sigma'^2_{\gamma} \approx \sigma^2_{\gamma}$), the new CV is:
\begin{equation}
    \CV' = \frac{\sigma_{\gamma}}{\gamma - |\Delta\gamma|} > \frac{\sigma_{\gamma}}{\gamma} = \CV
\end{equation}
since $\gamma - |\Delta\gamma| < \gamma$. The increase is:
\begin{equation}
    \Delta\CV = \CV' - \CV = \sigma_{\gamma} \left(\frac{1}{\gamma - |\Delta\gamma|} - \frac{1}{\gamma}\right) = \frac{\sigma_{\gamma} \cdot |\Delta\gamma|}{\gamma(\gamma - |\Delta\gamma|)}
\end{equation}
This is strictly positive for $|\Delta\gamma| < \gamma$, which holds for all methods tested (max reduction: 37\% < 100\%). The CNR bound follows from rearranging:
\begin{equation}
    \text{CNR} = \frac{|\Delta\gamma|}{\Delta\CV} = \frac{\gamma(\gamma - |\Delta\gamma|)}{\sigma_{\gamma}} \leq \frac{\gamma^2}{\sigma_{\gamma}} = \frac{\bar{\gamma}}{\CV}
\end{equation}
For the empirical setting ($\bar{\gamma} = 0.991$, $\CV = 0.157$), this gives $\text{CNR} \leq 6.31$.
\end{proof}

\begin{corollary}[Tradeoff Irreducibility]
The coupling-noise tradeoff is irreducible under the conditions of Theorem~\ref{thm:tradeoff}; outside those conditions the conclusion need not hold. In particular, the companion manuscript~\citep{companion2026} reports a 30-seed confidence-gated result ($\Delta\gamma = -0.332$ with $\Delta\CV < 0$), and our deterministic-seed simulation shows negative mean $\Delta\CV$ for all five methods ($-3.5\%$). These results are consistent with each other but violate the unqualified tradeoff claim, implying the constant-variance assumption in the proof is violated in these regimes. We therefore treat the tradeoff as conditional rather than universal; the completed joint-measurement experiments (Section~\ref{sec:joint}) found no significant coupling reduction for either factor at either horizon, consistent with this scoping.
\end{corollary}

\begin{proof}
Immediate from Theorem~\ref{thm:tradeoff}: if $\Delta\gamma < 0$, then $\Delta\CV > 0$ under the stated conditions.
\end{proof}

% ============================================================
% \S3 RESULTS
% ============================================================
\section{Results}

\subsection{Finding 1: The Claim Under Test --- Reported Temporal Calibration Fatigue}

\textbf{Framing.} This subsection reports the claim that motivated the study, as reported from the original GPT-4o protocol runs. Its per-repetition traces are not archived in this package and the values below cannot be independently recomputed (Limitations~(i)); we therefore treat them as \emph{the claim under test}, and report our replication attempts at the end of this subsection.

\textbf{Statistical significance summary.} All pairwise comparisons among the eight GPT-4o conditions in this section use Welch's $t$-test with Bonferroni correction for multiple comparisons; the qwen2.5:32b replication reported below instead uses paired (same-seed) $t$-tests, which is the appropriate pairing for its seed-matched design. Effect sizes are reported as Cohen's $d$ with 95\% bootstrap confidence intervals ($B{=}5000$). The key contrasts are: (i) C1 vs.~C1b (pure temporal): $t(58) = 12.4$, $p < 0.001$, $d = 4.85$ [3.82, 5.98]; (ii) C1b vs.~C2 (pure communication): $t(58) = 2.1$, $p = 0.04$, $d = 0.54$ [0.03, 1.05]; (iii) C1 vs.~C2 (total effect): $t(58) = 11.8$, $p < 0.001$, $d = 4.70$ [3.68, 5.81]. The temporal effect is approximately 9 times larger than the communication effect in standardized units. All statistics in this paragraph are reported from the GPT-4o runs (per-rep traces not archived; see Limitations).

Table~\ref{tab:results} presents the full eight-condition results. The dominant pattern is \textit{vertical stratification}: conditions at $T{=}10$ cluster at ECE $\approx 0.01$--$0.02$, while conditions at $T{=}20$ cluster at ECE $\approx 0.21$--$0.24$. Within each $T$ band, synchronization makes minimal difference.

\begin{table}[H]
\centering
\caption{\textbf{Full results matrix with isolation baselines.} C1b and C3b enable pure factor decomposition. Final ECE: means with 95\% bootstrap CIs ($B{=}5{,}000$), reported from the GPT-4o runs (per-rep traces not archived; see Limitations).}
\label{tab:results}
\begin{tabular}{@{}lccccc@{}}
\toprule
Cond. & $N$ & $T$ & Sync & Overall ECE & Final ECE [95\% CI] \\
\midrule
C1  & 3 & 10 & No  & $0.297$ & $0.012$ [$0.006$, $0.022$] \\
Q5  & 3 & 10 & Yes & $0.309$ & $0.035$ [$0.025$, $0.047$] \\
C3b & 5 & 10 & No  & $0.295$ & $0.017$ [$0.012$, $0.022$] \\
C3  & 5 & 10 & Yes & $0.293$ & $0.018$ [$0.014$, $0.024$] \\
\midrule
Q2  & 3 & 20 & No$^\dagger$ & $0.324$ & $0.208$ [$0.189$, $0.227$] \\
C1b & 3 & 20 & No  & $0.321$ & $0.219$ [$0.198$, $0.242$] \\
C2  & 3 & 20 & Yes & $0.325$ & $0.241$ [$0.219$, $0.263$] \\
C4  & 5 & 20 & Yes & $0.317$ & $0.234$ [$0.222$, $0.246$] \\
\bottomrule
\end{tabular}
\smallskip
{\footnotesize $^\dagger$Q2: no-history condition---agents do not see their own prior rounds.}
\end{table}

\begin{figure}[H]
\centering
\includegraphics[width=0.95\linewidth]{fig1_ece_trajectories.pdf}
\caption{Historical GPT-4o claim under test (not replicated): calibration fatigue trajectories across conditions. ECE trajectories for four conditions: isolated baseline C1 ($N{=}3$, $T{=}10$; green), communicating long-horizon C2 ($N{=}3$, $T{=}20$; blue) and C4 ($N{=}5$, $T{=}20$; red), and communicating short-horizon C3 ($N{=}5$, $T{=}10$; orange). Long-horizon sync conditions reach high ECE (C2 $0.241$, C4 $0.234$) while short-horizon conditions remain calibrated (C1 $0.012$, C3 $0.018$); the pure-factor decomposition (90.8\% temporal) is in Table~\ref{tab:results}. Shaded regions denote 95\% bootstrap confidence intervals ($B{=}5000$).}
\label{fig:ece_trajectories}
\end{figure}


\begin{table}[H]
\centering
\caption{\textbf{Pure factor decomposition (historical GPT-4o claim under test).} Temporal horizon accounts for 90.8\% of the original ``contagion'' effect (statistics reported from the GPT-4o runs; per-rep traces not archived, see Limitations).}
\label{tab:decomp}
\begin{tabular}{@{}lccccc@{}}
\toprule
Effect & Contrast & $\Delta$ECE [95\% CI] & \% Total & $p$ & $d$ \\
\midrule
Total (original) & C2 $-$ C1 & $+0.228$ [$+0.205$, $+0.253$] & 100\% & $<0.0001$ & 4.70 \\
Pure $T$ (no sync) & C1b $-$ C1 & $+0.207$ [$+0.181$, $+0.234$] & 90.8\% & $<0.0001$ & 4.85 \\
Pure Sync ($T{=}20$) & C2 $-$ C1b & $+0.021$ [$+0.003$, $+0.041$] & 9.2\% & $0.041$ & 0.54 \\
Pure $N$ (no sync) & C3b $-$ C1 & $+0.004$ [$+0.000$, $+0.012$] & 1.8\% & $0.284$ & 0.22 \\
Pure Sync ($T{=}10$) & C3 $-$ C3b & $+0.002$ [$-0.003$, $+0.008$] & 0.9\% & $0.761$ & 0.06 \\
\bottomrule
\end{tabular}
\end{table}

\textbf{Key results:} (i) Temporal horizon alone accounts for 90.8\% of the total effect ($\Delta_{\text{pure } T} = +0.207$, $d = 4.85$). (ii) Synchronization adds at most $+0.021$ (9.2\%) at $T{=}20$ (raw $p = 0.041$), with negligible effect at $T{=}10$ ($+0.002$); the pure-sync contrast is not significant after family-level Bonferroni over the five decomposition contrasts ($k{=}5$: adjusted $p = 0.205$), so the 9.2\% attribution is suggestive. (iii) Agent count is negligible ($+0.004$, $p = 0.284$). (iv) Context accumulation is ruled out: Q2 (no-history, ECE $= 0.208$) overlaps with C1b (with-history, ECE $= 0.219$; 95\% CIs overlap).

\textbf{Accuracy trajectories.} Accuracy declines modestly ($\approx 0.69$ at $T{=}10$ to $\approx 0.65$ at $T{=}20$), but ECE increases disproportionately ($+0.21$), confirming calibration fatigue is an alignment problem, not accuracy collapse~\citep{dodge2019show}.

\textbf{Cross-model validation.} Table~\ref{tab:cross} extends to DeepSeek V4 Pro and V4 Flash. DeepSeek models show minimal temporal fatigue ($\Delta$ECE $< 0.06$ at $T{=}20$), confirming the overconfident seed agent is necessary for synchronization effects. Without a seed agent (R2), communication does not increase ECE.

\begin{table}[H]
\centering
\caption{\textbf{Cross-model validation.} DeepSeek models are naturally well-calibrated and resist temporal fatigue.}
\label{tab:cross}
\begin{tabular}{@{}lcccc@{}}
\toprule
Model & C1@10 (ISO) & C4 (SYNC) & R1@20 (ISO) & R2 (SYNC, no seed) \\
\midrule
GPT-4o & $0.012$ & $0.234$ & $0.219$ & --- \\
V4 Pro & --- & --- & $0.088$ & $0.072$ \\
V4 Flash & --- & --- & $0.031$ & $0.030$ \\
\bottomrule
\end{tabular}
\end{table}

\paragraph{Reproduction status.} The eight-condition decomposition above is reported from the original GPT-4o protocol runs. An independent reproduction attempt on deepseek-v4-flash (trap-question causal control, three repetitions per condition; \texttt{evidence/p1\_causal\_quick\_20260808.json}; calibration pressure probe \texttt{evidence/p1\_calibration\_pressure\_probe\_20260808.json}) did not reproduce the temporal fatigue effect: C2 final ECE at $T{=}20$ overlapped the $T{=}10$ range, and the no-peer/placebo control ordering was inconsistent with the sync-driven mechanism. The $+0.207$/90.8\% decomposition therefore appears GPT-4o-specific. A second, full-protocol independent replication on qwen2.5:32b (all eight conditions of the design matrix, 30 paired seeds per condition, temperature 0, seed-locked; 240/240 runs with all per-round records archived in \texttt{evidence/finding1\_rerun\_qwen32b/}) reproduced the total C2$-$C1 direction ($+\Delta$ECE $= +0.078$, $p = 0.002$, $d = 0.81$) but with a \emph{reversed effect structure}: the pure temporal contrast was statistically indistinguishable from zero ($-0.004$, $p = 0.86$, $d = -0.04$) while the synchronization contrast---a minor term (9.2\%) in the GPT-4o decomposition---became the dominant significant terms (T20 sync: $+0.082$, $p < 0.0001$, $d = 0.85$; T10 sync: $+0.079$, $p = 7.9{\times}10^{-8}$, $d = 1.25$; pure $N$ remained null). The temporal-fatigue interpretation is thus GPT-4o-specific across two independent model families; the 90.8\% attribution must not be read as a model-general claim. As we show next, the communicative effect is not model-general either. Absolute ECE levels differ substantially between the rerun (condition means $\approx 0.59$--$0.67$) and Table~\ref{tab:results} ($\approx 0.01$--$0.24$); only contrast directions, not levels, are compared here.

\textbf{A third full-protocol replication on a newer backbone.} We repeated the complete eight-condition protocol on Qwen3.8-27B (AWQ-INT4 served via vLLM with thinking disabled; 30 paired seeds per condition, 240/240 runs, all per-round records archived under \texttt{evidence/finding1\_rerun\_qwen38\_27b/}; condition means $0.60$--$0.63$, comparable to the qwen2.5:32b rerun, so the two are protocol-comparable). Again the pure temporal contrast is statistically indistinguishable from zero ($+0.0003$, $p = 0.99$, $d = 0.002$). Unlike qwen2.5:32b, however, neither synchronization contrast reaches significance here (T20: $+0.029$, $p = 0.21$, $d = 0.23$; T10: $+0.022$, $p = 0.12$, $d = 0.30$) and the total effect is null ($+0.029$, $p = 0.42$). We therefore do not describe the communicative effect as cross-model robust either: what replicates across all three models is the \emph{absence} of temporal fatigue, not the presence of a communication effect.

\begin{figure}[t]
\centering
\includegraphics[width=\linewidth]{figures/fig9_crossmodel_forest.pdf}
\caption{\textbf{The temporal contrast does not replicate, and the communication contrast is model-dependent.} Left: the pure temporal contrast ($C1b-C1$) with 95\% bootstrap CIs over 30 paired seeds for both full-protocol replications (filled circles) and the reported GPT-4o value (open square; traces not archived). Both replication CIs straddle zero \emph{and exclude} the reported $+0.207$. Right: the communication contrast ($C2-C1b$), significant in qwen2.5:32b but not in Qwen3.8-27B, so this component is not model-general.}
\label{fig:crossmodel}
\end{figure}

Figure~\ref{fig:crossmodel} summarises both replications: the temporal contrast is zero everywhere and its confidence intervals exclude the reported effect size, while the communication contrast is significant in one model and null in the other.

\textbf{Why these effects are hard to detect: the instrument itself.} Our archive also indicates why such contrasts are fragile. Across the 26{,}400 per-round records of the two full replications, self-reported confidence separates correct from incorrect answers by only $+0.015$ (qwen2.5:32b) and $+0.010$ (Qwen3.8-27B), and takes just 9 and 5 distinct values respectively. Near-zero \emph{answer-dependence} of verbalized confidence has been identified as a primary driver of LLM overconfidence by \citet{advice2026}; our multi-model archive confirms that it holds across model families and helps explain the pattern we report: because reported confidence barely varies with correctness, ECE in these experiments is dominated by the confidence \emph{scale}, leaving little dynamic range in which a temporal or communicative effect could appear. This is consistent with the exactly-zero temporal contrasts above and with the model-dependence of the communication contrast, and it bounds what any study using verbalized confidence as its measurement instrument can detect.

\textbf{A smaller-scale pilot (low power).} A methodologically independent local rebuild of the same condition structure on a smaller model (qwen2.5:7b; $k{=}5$ agents, 30-question pool, three repetitions per condition; script and raw output archived under \texttt{evidence/finding1\_rerun\_qwen25\_7b\_pilot/}) shows the same qualitative signature: the isolated horizon contrast is $\approx 0$ ($-0.004$) while synchronization effects are large ($+0.081$ to $+0.127$). We report this as a \emph{pilot}, not a replication: with three repetitions it cannot resolve small effects, it uses a reduced agent/round configuration, it is our own rebuild rather than a third-party implementation, and its absolute ECE scale ($0.03$--$0.24$) matches neither the eight-condition protocol nor the 32b rerun, so it is not directly comparable to either and it carries no weight in any pooled statistic.
\subsection{Finding 2: Self-Evaluation Immunity Is Model-Dependent (Descriptive)}

Prior work~\citep{epc2026} reported DeepSeek-chat self-evaluation produced $\gamma \approx 0.03$ ($N{=}30$, 97\% zero coupling). We replicate on DeepSeek-V4-Pro ($N{=}5$, within-subjects; per-seed API traces are not archived, so the values below are descriptive rather than independently recomputable---see Limitations~(ii)).

\begin{table}[H]
\centering
\caption{\textbf{Self-evaluation on DS-V4-Pro.} Standard self-eval produces coupling ($\gamma \approx 1.06$), contradicting prior immunity (per-seed API traces not archived; see Limitations).}
\begin{tabular}{@{}lcccc@{}}
\toprule
Mode & $\bar{\gamma}_{T \to V}$ & $\bar{\gamma}_{V \to T}$ & $\text{JSD}_{T \to V}$ & $\text{JSD}_{V \to T}$ \\
\midrule
Standard self-eval & $1.064$ & $0.988$ & $0.213$ & $0.191$ \\
Calibrated self-eval & $\mathbf{0.617}$ & $\mathbf{0.752}$ & $\mathbf{0.071}$ & $\mathbf{0.096}$ \\
\midrule
$\Delta$ (reduction) & $-42\%$ & $-24\%$ & $-67\%$ & $-50\%$ \\
\bottomrule
\multicolumn{5}{l}{\footnotesize Paired bootstrap 95\% CIs ($B{=}2000$): $\gamma_{T\to V}$ [0.76, 1.41] (std) vs.\ [0.31, 0.98] (cal).} \\
\end{tabular}
\end{table}

\begin{figure}[H]
\centering
\includegraphics[width=0.9\linewidth]{fig2_final_ece.pdf}
\caption{Final ECE comparison across all eight conditions (reported GPT-4o runs; per-rep traces not archived, see Limitations). Shared short-horizon calibration holds ECE at 0.012 in C1; extending the horizon reintroduces miscalibration in the $T{=}20$ sweep conditions (C1b 0.219, C2 0.241) while synchronization adds only $+0.021$. Within these reported runs temporal horizon---not communication---accounts for the dominant share; our independent replications do not reproduce this attribution. Error bars denote 95\% bootstrap CIs.}
\label{fig:final_ece}
\end{figure}


Self-evaluation immunity does not generalize across model versions. DS-V4-Pro produces $\gamma_{T \to V} = 1.064$---comparable to cross-model coupling---directly contradicting prior near-zero results~\citep{epc2026}. Calibration directionally reduces coupling ($-42\%$)~\citep{tian2023just}. The safest evaluator configuration is not guaranteed safe; model updates can silently invert calibration behavior~\citep{sclar2024quantifying}.

\subsection{Finding 3: Coupling Reduction without a Robust Noise Increase in Simulation}

\textbf{3a: Noise (negative result).} We test five calibration methods across 60 simulation conditions (5 methods $\times$ 4 $\alpha_{\text{eval}}$ levels $\{0.0, 0.5, 1.0, 2.0\}$ $\times$ 3 $\tau$ levels $\{0.5, 1.0, 2.0\}$, excluding the no-calibration baseline; canonical script with a deterministic per-method seed, stable across five independent runs). Table~\ref{tab:noise} shows that none of the five methods increases mean CV in the $\alpha{>}0$ subset: mean $\Delta\CV = -3.5\%$ (range $-7.0\%$ to $-0.6\%$), with per-method CV$\downarrow$rate between 56\% and 78\% (most conditions show decreased CV). Earlier script versions used Python's hash-based seeds and produced sign-varying $\Delta\CV$ across runs, so the simulation grid does not robustly support a universal noise-increase claim; the coupling--noise tradeoff remains an open hypothesis for the pre-specified joint experiment (\S\ref{sec:joint}). $\Delta$ECE and $\Delta$CV are uncorrelated ($r = -0.04$, $p = 0.95$, $n{=}5$), a pattern reminiscent of calibration--reliability tensions reported for elicited confidence~\citep{tian2023just}; whether a general ECE--CV dissociation holds beyond that setting remains to be established.

\begin{table}[H]
\centering
\caption{\textbf{Calibration noise effects by method (canonical script, deterministic seed, stable across 5 runs; $\alpha{>}0$, $n{=}9$).} None of the five methods increases mean CV in this grid; $\Delta\CV$ is seed-sensitive across script versions. $p$: two-sided binomial test vs.\ 50\% on CV$\downarrow$rate; none significant after Bonferroni ($\alpha' = 0.010$).}
\label{tab:noise}
\begin{tabular}{@{}lccccc@{}}
\toprule
Method & $\Delta$ECE & CV$\downarrow$rate & $p$ & $\Delta$CV & CV SE \\
\midrule
Platt & $+0.0131$ & 5/9 (56\%) & 1.00 & $-2.2\%$ & $\pm 0.004$ \\
Temperature & $-0.0038$ & 6/9 (67\%) & 0.51 & $-6.5\%$ & $\pm 0.004$ \\
Isotonic & $-0.0226$ & 6/9 (67\%) & 0.51 & $-1.2\%$ & $\pm 0.004$ \\
Histogram & $-0.0219$ & 6/9 (67\%) & 0.51 & $-0.6\%$ & $\pm 0.004$ \\
\midrule
\textbf{Conf-Gating} & $\mathbf{-0.0226}$ & 7/9 (78\%) & 0.18 & $\mathbf{-7.0\%}$ & $\pm 0.004$ \\
\bottomrule
\multicolumn{6}{l}{\footnotesize $p$: two-sided binomial test vs.\ 50\%. CV SE retained from the original grid; conf-gating uses additive update + isotonic calibration.} \\
\end{tabular}
\end{table}


\textbf{Robustness and reproducibility.} The canonical simulation is fully deterministic (per-method seed offsets replacing the earlier hash-based seeds); five independent runs produced byte-identical outputs (SHA-256 \texttt{FB0D7815B805\ldots}). Per-method bootstrap 95\% CIs for mean $\Delta\CV$ ($n{=}9$ conditions, $B{=}2000$, seed 42) all include zero: Platt $-2.2\%$ [$-8.7, +4.4$], Temperature $-6.5\%$ [$-14.4, +0.7$], Isotonic $-1.2\%$ [$-11.6, +10.8$], Histogram $-0.6\%$ [$-6.5, +5.7$], Conf-Gating $-7.0\%$ [$-14.8, +1.7$]. None of the five methods shows a statistically significant noise increase---or decrease---in this grid, and the sign-varying $\Delta\CV$ of earlier hash-seeded script versions confirms that the noise effect is condition- and seed-sensitive. We therefore do not claim a universal noise increase; the coupling--noise tradeoff is addressed by the completed joint experiment reported in \S\ref{sec:joint}. A synthetic second-evaluator-profile check (Qwen-3.7-like evaluator preferences, same executor task profiles; \texttt{evidence/finding3\_second\_domain.json}) gives mean $\Delta\CV = +5.7\%$ (95\% CI $[-13.7, +31.3]$, $n{=}9$ conditions, 6/9 negative): not a significant increase, but the per-condition direction is profile-dependent, so the no-noise-increase conclusion is specific to the primary evaluator.
\textbf{3b: Coupling reduction.} DeepSeek-V4-Pro executor, GLM5.2 evaluator, within-subjects ($N{=}15$, $R{=}30$). Table~\ref{tab:coupling} shows confidence-gated TTRL significantly reduces coupling with non-overlapping bootstrap CIs.

\begin{table}[H]
\centering
\caption{\textbf{Confidence-gated TTRL reduces coupling.} Bootstrap 95\% CIs ($B{=}2000$) do not overlap between modes.}
\label{tab:coupling}
\begin{tabular}{@{}lccccc@{}}
\toprule
Mode & $\bar{\gamma}_{T \to V}$ [95\% CI] & $\bar{\gamma}_{V \to T}$ [95\% CI] & JSD$_{T \to V}$ & JSD$_{V \to T}$ \\
\midrule
Standard TTRL & $0.991$ [$0.910$, $1.086$] & $0.938$ [$0.841$, $1.037$] & $0.195$ & $0.173$ \\
Calibrated TTRL & $0.623$ [$0.556$, $0.682$] & $0.606$ [$0.555$, $0.657$] & $0.071$ & $0.069$ \\
\midrule
$\Delta$ (reduction) & $-37.2\%$ & $-35.4\%$ & $-63.4\%$ & $-60.2\%$ \\
\bottomrule
\end{tabular}
\end{table}

% ============================================================
% \S4 CALIBRATION LANDSCAPE
% ============================================================
\section{The Calibration Landscape}

\subsection{Unified Framework}

Synthesizing Findings~1--3 with the 11-condition evaluator survey, we define the \textbf{calibration landscape}: a four-quadrant mapping of evaluator configurations.

\begin{table}[H]
\centering
\caption{\textbf{The calibration landscape quadrant map.} The favorable regime is empty. The final column states each entry's \emph{evidentiary status} rather than its importance: ``Unverifiable'' marks a value whose per-repetition traces are not archived, ``Descriptive'' a value reported without archived traces, and ``Verified'' a result recomputable from the shipped package.}
\begin{tabular}{@{}lllccc@{}}
\toprule
Quadrant & Q & Experiment & $N$ & Key Result & Evidence \\
\midrule
Target: Cost & Q1 & Temporal fatigue & 30 & ECE $\uparrow$, $d{=}4.85$ & Unverifiable (traces absent) \\
Target: Cost & Q2 & Self-eval fragility & 5/30 & Immunity version-dependent & Descriptive (traces absent) \\
Tool: Cost & Q3a & Noise increase & 60 (sim) & Mean $\Delta$CV $\approx -3.5$\% (range $-7.0$ to $-0.6$) & Verified (simulation) \\
Tool: Benefit & Q3b & Coupling reduction & 15 & $\gamma \downarrow 35$--$37\%$ & Descriptive (traces absent) \\
\bottomrule
\end{tabular}
\end{table}

Three quadrants identify costs/fragilities; one identifies a benefit. The landscape is asymmetric: calibration methods that reduce coupling do not robustly increase noise in the simulation grid (mean $\Delta\CV = -3.5\%$ under the deterministic seed).

\textbf{On the two coupling-reduction figures.} The $-42\%$ coupling reduction of Finding~2 (self-evaluation, within-subjects, $N{=}5$) and the $\gamma \downarrow 35$--$37\%$ benefit above (confidence-gated TTRL, $N{=}15$, $R{=}30$) come from \emph{different designs} and are not two estimates of one quantity. Both are descriptive: neither run's per-seed traces are archived in this package (Limitations~(ii)).

\subsection{Formal Hypothesis}

\begin{hypothesis}[Calibration Tradeoff]\label{hyp:tradeoff}
In TTRL-based agent adaptation, calibration methods that reduce evaluator preference coupling ($\Delta\gamma < 0$) will simultaneously increase measurement noise ($\Delta\CV > 0$).
\end{hypothesis}

\textbf{Falsification condition:} A method satisfying both $\Delta\gamma < 0$ (significant) and $\Delta\CV \leq 0$ in a single joint experiment with $N \geq 30$.

% ============================================================
\subsection{Joint-Measurement Protocol: Completed Experiment}\label{sec:joint}

\subsubsection{Design and Data Quality}

We contribute and execute a protocol for testing Hypothesis~\ref{hyp:tradeoff}. Design: $2 \times 2$ factorial (Communication on/off $\times$ Calibration on/off), confidence-gated TTRL, $N{=}30$ seeds per condition ($120$ total), $T{=}30$ rounds, $K{=}5$ agents, simultaneously recording $\gamma$, $H$, $\CV(N)$, and ECE across all seeds. The four conditions are: (A) No Communication, No Calibration (baseline); (B) Communication, No Calibration; (C) No Communication, Calibration; (D) Communication + Calibration. The experiment was executed on Qwen3-14B-AWQ via vLLM (cloud RTX 3090), requiring approximately $54{,}000$ API calls. Per-seed data and summaries are archived under \texttt{evidence/f9\_uniform\_T30\_results/} (\texttt{analysis\_summary.json}, per-seed JSONs, \texttt{joint\_experiment\_T30.py}); the T50 extension per-seed data under \texttt{evidence/f9\_t50\_results/} (120 seeds, 30/condition) with analysis in \texttt{evidence/f9\_T50\_analysis/analysis\_summary.json} (\texttt{evidence/f9\_T50\_full\_analysis.py}).

\textbf{Data quality.} All $120$ seeds passed automated quality checks: ECE $\in [0.075, 0.369]$ (all below the $0.40$ suspicious-data threshold), accuracy $\in [0.587, 0.887]$ (all above the $0.30$ floor), and no anomalous clustering at the default confidence level. Zero seeds were excluded.

\subsubsection{Overall Result: No Effect at Convergence}

A $2 \times 2$ ANOVA on the 30-round average reveals no significant main effects or interactions for any metric (all $p > 0.12$, smallest $p = 0.123$; Cohen's $d < 0.29$). Table~\ref{tab:joint_overall} reports the descriptive statistics. The multi-agent system is robust to both Communication and Calibration manipulations at convergence: final-state metrics are nearly identical across conditions. Table~\ref{tab:joint_overall} and the ANOVA values above derive from the 120 per-seed JSONs as the single canonical source (regenerated 2026-09-10 via \texttt{evidence/regen\_T30\_singlesource\_20260910.py} after 24/120 seeds were found to disagree with a prior CSV build; main effects Student's $t$, $d$ with population-SD pooling, interaction via $2{\times}2$ ANOVA $F$).

\textbf{Power of this null.} With $n{=}30$ seeds per cell (and $60$ vs.\ $60$ for each main effect), the 80\%-power minimum detectable effect is ${\approx}0.72$~SD for a cell-to-cell contrast and ${\approx}0.51$~SD for a main effect (two-sample, $\alpha{=}0.05$, pooled SD from the archived per-repetition records). All effects reported in this subsection are $|d|{<}0.29$, well below that threshold, so the null is informative rather than merely underpowered.

\begin{table}[H]
\centering
\caption{\textbf{Joint-measurement experiment: 30-round average descriptive statistics} ($N{=}30$ seeds per condition, mean $\pm$ SD). Lower ECE and $H$ are better; higher $\gamma$ and accuracy are better.}
\label{tab:joint_overall}
\begin{tabular}{@{}lcccc@{}}
\toprule
Metric & A (Baseline) & B (Comm) & C (Cal) & D (Comm+Cal) \\
\midrule
ECE $\downarrow$ & $0.229 \pm 0.070$ & $0.226 \pm 0.072$ & $0.232 \pm 0.065$ & $0.225 \pm 0.071$ \\
$\gamma \uparrow$ & $1.147 \pm 0.142$ & $1.113 \pm 0.139$ & $1.151 \pm 0.136$ & $1.128 \pm 0.122$ \\
$H \downarrow$ & $0.909 \pm 0.034$ & $0.918 \pm 0.027$ & $0.909 \pm 0.038$ & $0.918 \pm 0.028$ \\
Accuracy $\uparrow$ & $0.723 \pm 0.070$ & $0.727 \pm 0.073$ & $0.719 \pm 0.066$ & $0.728 \pm 0.073$ \\
\bottomrule
\end{tabular}
\end{table}

The full ANOVA table confirms the null result: Communication main effect on ECE ($d{=}-0.072$, $p{=}0.696$), $\gamma$ ($d{=}-0.213$, $p{=}0.251$), $H$ ($d{=}+0.286$, $p{=}0.123$); Calibration main effect on all metrics ($|d| < 0.08$, all $p > 0.70$); Communication $\times$ Calibration interaction (all $p > 0.82$). The joint null is complete for ECE (interaction $p=0.91$), but per-round $\gamma$ shows a descriptive early-positive to late-negative drift, and $H$ shows a weak comm trend ($d=0.29$, $p=0.12$); neither supports a claim at $N=120$.

\subsubsection{Late-Phase Communication Pattern: Suggestive but Not Robust}

While the 30-round average shows no effect, per-round analysis suggests a \textit{delayed} Communication pattern that, if real, emerges only in the late phase. We partition the 30 rounds into early (Rounds 1--5) and late (Rounds 26--30) windows and test each separately. Throughout, the individual seed is the unit of analysis (one window-averaged observation per seed, Welch tests). An earlier version of this analysis pooled the five within-seed rounds as if they were independent observations; this pseudo-replication inflates significance by construction and produced the $p {=} 0.009$--$0.015$ reported in a previous draft. All values below are seed-level.

\textbf{Early phase (Rounds 1--5):} No significant effects. Communication on $\gamma$: $d{=}+0.096$, $p{=}0.600$; on $H$: $d{=}-0.086$, $p{=}0.638$. Calibration and interaction are also null.

\textbf{Late phase (Rounds 26--30):} The Communication pattern is directionally consistent with a delayed effect, but does not reach significance at the seed level:

\begin{itemize}[leftmargin=*,nosep]
    \item Communication \textit{reduces} consensus strength $\gamma$: $d{=}-0.245$, $p{=}0.182$ (late-phase means: NoComm $=0.637$, Comm $=0.605$)
    \item Communication \textit{increases} strategy entropy $H$: $d{=}+0.309$, $p{=}0.093$ (late-phase means: NoComm $=0.916$, Comm $=0.925$)
    \item Calibration shows no significant effect at any phase (all $p > 0.38$)
    \item No significant Communication $\times$ Calibration interaction at any phase
\end{itemize}

Table~\ref{tab:joint_early_late} summarizes the early-versus-late contrast. The pattern is not present initially and is strongest as the system converges, consistent with a \textit{time-delayed} influence---but the evidence is a non-significant trend, not an established effect.

\begin{table}[H]
\centering
\caption{\textbf{Early vs.\ late phase Communication effect} (Cohen's $d$ and seed-level Welch $p$ on 5-round window means; one observation per seed, $n{=}60$ per merged arm). Significant effects ($p < 0.05$) in bold; none survive.}
\label{tab:joint_early_late}
\begin{tabular}{@{}lcccc@{}}
\toprule
 & \multicolumn{2}{c}{Early (R1--5)} & \multicolumn{2}{c}{Late (R26--30)} \\
\cmidrule(lr){2-3} \cmidrule(lr){4-5}
Metric & Comm $d$ & $p$ & Comm $d$ & $p$ \\
\midrule
$\gamma$ & $+0.096$ & $0.600$ & $-0.245$ & $0.182$ \\
$H$ & $-0.086$ & $0.638$ & $+0.309$ & $0.093$ \\
\bottomrule
\end{tabular}
\end{table}

A per-round analysis at $T{=}30$ is compatible with a growing trend---single-round Communication effects on $\gamma$ and $H$ are largest in the mid-to-late phase ($|d| \approx 0.24$--$0.28$)---but no single round reaches significance at the seed level (all $p \geq 0.17$), so the pattern is a weak, non-significant trend rather than an established delayed effect. The $T{=}50$ extension below was designed to test exactly this trend.

\subsubsection{Mechanism: Slower Convergence and Maintained Diversity}

We investigate the mechanism that would underlie the suggestive late-phase Communication pattern through three analyses of agent strategy weights.

\textbf{Weight convergence speed.} Comm conditions show $7$--$8\%$ lower late-phase weight convergence (measured as $1 - H/H_{\max}$): A $=0.0838$, B $=0.0779$, C $=0.0831$, D $=0.0767$. Communication slows the rate at which agents concentrate weight on dominant strategies.

\textbf{Inter-agent strategy heterogeneity.} Comm conditions show $2.7\%$ higher pairwise Jensen--Shannon divergence (JSD) between agents' final weight vectors: A $=0.2105 \pm 0.028$, B $=0.2138 \pm 0.039$, C $=0.2083 \pm 0.035$, D $=0.2163 \pm 0.029$ (Communication effect: $d{=}+0.170$, $p{=}0.35$; direction consistent but not independently significant). Agents in Comm conditions maintain more diverse strategy distributions from each other.

\textbf{Strategy preference distribution.} In the no-Comm conditions, the \textit{algebraic} strategy dominates with weight $0.187$; in Comm conditions, this drops to $0.161$, and the overall distribution is more uniform across the strategy space. Communication weakens the concentration of strategy weight on any single approach.

The mechanism chain is: Communication $\rightarrow$ more uniform strategy exploration $\rightarrow$ slower weight convergence $\rightarrow$ higher inter-agent heterogeneity $\rightarrow$ lower late-phase $\gamma$, higher $H$. At $T{=}30$ this chain is consistent with the descriptives above; the $T{=}50$ extension below finds the same directions but none of these components reaches significance, so we treat the chain as a hypothesis rather than a demonstrated mechanism.

\subsubsection{$T{=}50$ Horizon Extension: Pre-specified Strengthening Test}

The $T{=}30$ analysis framed the late-phase pattern as a potentially cumulative effect and pre-specified a falsification criterion: extending to $T{=}50$ rounds under the identical $2 \times 2$ design, a genuinely cumulative Communication effect should reach $d > 0.30$ in the Rounds 46--50 window. We completed this extension ($N{=}120$ seeds, same Qwen3-14B-AWQ via vLLM setup; all seeds passed the same quality checks: ECE $\in [0.129, 0.374]$, accuracy $\in [0.580, 0.824]$, zero exclusions; ranges verified by per-seed recomputation over \texttt{evidence/f9\_t50\_results/} --- the \texttt{analysis\_summary.json} overall-ECE block is NaN and is not the source of these ranges). The criterion was \textbf{not met}: the Rounds 46--50 Communication effect is $d{=}-0.169$ ($p{=}0.356$) on $\gamma$ and $d{=}+0.208$ ($p{=}0.256$) on $H$. Table~\ref{tab:t50_extension} summarizes the seed-level window tests in both runs.

\begin{table}[H]
\centering
\caption{\textbf{Communication effect across horizons} (seed-level Cohen's $d$, Welch $p$; one observation per seed). The pre-specified strengthening criterion ($d > 0.30$ in R46--50) is not met; no window reaches significance in either run, but the sign of every window effect replicates across horizons.}
\label{tab:t50_extension}
\begin{tabular}{@{}llcc@{}}
\toprule
Window & Metric & $T{=}30$ & $T{=}50$ \\
\midrule
R26--30 & $\gamma$ & $-0.245$ $(p{=}0.182)$ & $-0.123$ $(p{=}0.502)$ \\
R26--30 & $H$ & $+0.309$ $(p{=}0.093)$ & $+0.148$ $(p{=}0.419)$ \\
R46--50 & $\gamma$ & --- & $-0.169$ $(p{=}0.356)$ \\
R46--50 & $H$ & --- & $+0.208$ $(p{=}0.256)$ \\
\bottomrule
\end{tabular}
\end{table}

At $T{=}50$, the overall (50-round average) Communication effect on $\gamma$ is $d{=}-0.177$ ($p{=}0.334$), and the mechanism components that motivated the delayed-effect interpretation shrink to negligible: the convergence delay between Comm and NoComm conditions is $0.3$ rounds against a $\sim$$38$--$43$-round convergence point, and the final-$H$ effect is $d{=}+0.225$ ($p{=}0.221$). Figure~\ref{fig:t30_vs_t50} compares the window-level Communication effect on $\gamma$ across the two horizons: the R26--30 effect shrinks from $d{=}-0.245$ to $d{=}-0.123$ rather than amplifying with horizon, and the R46--50 window at $T{=}50$ remains null.

\begin{figure}[H]
\centering
\includegraphics[width=0.75\linewidth]{fig8_T30_vs_T50.pdf}
\caption{Window-level Communication effect on $\gamma$ (Cohen's $d$ $\pm$ 95\% CI, seed-level) at $T{=}30$ vs.\ $T{=}50$. The effect does not strengthen with the longer horizon; no window is significant at the seed level in either run.}
\label{fig:t30_vs_t50}
\end{figure}

We therefore report the delayed-Communication pattern as a \textit{fragile, non-significant trend} rather than an established finding: the sign pattern of every window effect replicates across horizons ($\gamma$: positive early, negative late; $H$: the reverse), which keeps a small-effect version of the hypothesis alive, but no seed-level test reaches significance at either horizon and the pre-specified strengthening criterion fails. The convergence-rate reframing of the coupling--noise tradeoff is correspondingly downgraded to speculation; within this $2 \times 2$, neither Communication nor Calibration robustly alters convergence dynamics at any horizon tested.

% ============================================================
% §7 PRACTICAL GUIDELINES
% ============================================================


\section{Related Work}

\textbf{Novelty positioning.} This work makes three distinct contributions relative to prior literature. \textbf{First}, prior work on calibration in multi-agent systems~\citep{huang2026cagecal,epc2026} reported a ``calibration contagion'' effect but did not deconfound temporal from communicative causes; we pre-specified and executed that deconfounding, and \emph{our independent replications fail to reproduce the reported temporal dominance} (90.8\% of the total mean-ECE difference in the unarchived GPT-4o runs), with the qwen2.5:32b rerun showing a reversed effect structure. \textbf{Second}, near-zero coupling ($\gamma \approx 0.03$) previously observed for self-evaluation immunity on DeepSeek-chat had been treated as a stable property; we observe that model updates can silently invert this behavior ($\gamma \approx 1.06$ on DeepSeek-V4-Pro; descriptive, traces unarchived), indicating that self-evaluation immunity is model-dependent and fragile. \textbf{Third}, prior work on the coupling-noise tradeoff~\citep{caleffects2026} proposed a structural tradeoff; our deterministic-seed rerun (mean $\Delta\CV = -3.5\%$) fails to support a universal noise-increase claim, and we reframe the tradeoff as a method-specific hypothesis with explicit falsification conditions (Theorem~\ref{thm:tradeoff}).

\textbf{Division of labor with the companion manuscript.} This paper is the falsification and replication core: the eight-condition GPT-4o claim under test, three independent replication attempts (deepseek-v4-flash causal control; qwen2.5:32b and Qwen3.8-27B full-protocol reruns, 240 runs each, all per-round records archived), the deterministic simulation grid, and the completed $2{\times}2$ joint experiments at $T{=}30$ and $T{=}50$. The companion manuscript~\citep{companion2026} carries the complementary 30-seed self-evaluation extension ($\Delta\gamma=-0.332$, $\Delta\CV<0$) and the conditional communication analysis; it shares the canonical evidence package but makes no replication claim about the 90.8\% decomposition. Readers should cite this manuscript for the negative replication and the joint-experiment nulls, and the companion for the extended self-evaluation sample. The two must not be pooled as independent evidence for any shared number. \textbf{Appendix division of labor (A-line).} Retrieval-side leakage (the tmlr-p19 line, supplied by the fixed-correctness cloud rerun), pilot-scale mitigation explorations (the tmlr-p9/p10 calibrated-TTRL and self-evaluation pilots, citable as protocol deposits only), and negative-result records (the frozen tmlr-p15 headline and the tmlr-p18 recalibration failure) are consolidated as appendices mapped in \texttt{MERGE\_MAP\_A.md}; the archived paper3 and merged-f1-preprint drafts are historical versions and are not submitted independently.

\textbf{LLM-as-judge and evaluator bias.} LLMs are widely used as evaluators in benchmark pipelines~\citep{zheng2023judging,chiang2024chatbot}. Documented biases include position bias, verbosity bias, and self-preference amplification. The EPC framework~\citep{epc2026} extended this line by studying how evaluator biases propagate through closed-loop agent adaptation, introducing diagnostic metrics ($\gamma$, JSD) and documenting coupling across 13 conditions. All prior work in this line stopped at diagnosis; this paper is the first to evaluate calibration as a mitigation strategy in closed-loop TTRL adaptation.

\textbf{Multi-agent LLM dynamics.} LLM agents in communicative settings exhibit convergence phenomena: representational collapse in ensembles, herding effects in reasoning, conformity in debate, and emergent consensus in annotation tasks. Several works have studied whether communication improves or harms collective outcomes---\citet{bertalanic2026cost} found isolated self-correction outperforms homogeneous debate, while \citet{huang2026cagecal} identified Communication-Induced Over-Confidence and proposed graph-based correction. Opinion dynamics models---DeGroot weighted averaging, Friedkin-Johnsen social influence, bounded confidence---predict cumulative convergence under repeated interaction. Our calibration contagion experiment (Finding~1) extends this literature by measuring calibration degradation rather than accuracy or consensus; in the reported GPT-4o runs the degradation stratifies by horizon (total $+0.228$; pure $N$ only $+0.004$, Table~\ref{tab:decomp}), a pattern our replications do not reproduce, so any temporal-accumulation reading is confined to those unarchived runs.

\textbf{Probability calibration.} Post-hoc calibration methods---Platt scaling, temperature scaling, isotonic regression, histogram binning---effectively correct miscalibration in classification~\citep{guo2017calibration,niculescu2005predicting}. Our work differs from prior calibration studies in setting (test-time TTRL without external reward) and objective (coupling reduction rather than accuracy improvement). Most critically, we measure calibration's \textit{side effects}---the ECE--CV dissociation (Finding~3) and calibration contagion (Finding~1)---which are invisible when calibration is evaluated only on static accuracy.

\textbf{Evaluation reliability.} The instability of LLM benchmark scores is well documented: \citet{dodge2019reproducibility} showed that standard deviations under resampling can exceed reported improvements, and \citet{sclar2024quantifying} characterized prompt sensitivity as a major source of variance. The TEE framework~\citep{messing2026tee} decomposes evaluation error and identifies multi-judge ensembles as the most effective variance-reduction strategy. The reliability paradox~\citep{messing2026tee} independently documents that calibration does not improve evaluation reliability. Our Finding~3 confirms this pattern and extends it to closed-loop TTRL evaluation with five methods including confidence-gating.

\textbf{Self-evaluation and model stability.} Self-evaluation is widely deployed in agent systems for its convenience and reported immunity to external bias. The finding that this immunity is model-version-dependent (Finding~2) adds to growing evidence of LLM evaluator instability: the GPT-4o version drift documented in~\citep{epc2026} showed that a silent API update can invert a study's qualitative conclusion. Our contribution is demonstrating that this instability extends to self-evaluation---the configuration previously considered most stable.

% ============================================================
% §9 DISCUSSION
% ============================================================


% \S5 DISCUSSION
% ============================================================

\subsection{Calibration in Multi-Agent Systems}

Recent work has begun to study calibration in multi-agent LLM settings. \citet{huang2026cagecal} document that communication can degrade calibration in debate-like settings, but their analysis does not decompose temporal from communicative effects. Our pre-specified deconfounding (Table~\ref{tab:decomp}) addresses this gap: the GPT-4o decomposition attributes 90.8\% of the observed degradation to temporal horizon alone, but our own replications do not reproduce that attribution, and the original per-repetition traces are not archived (Limitations~(i)). Prior reports that self-evaluation produces near-complete immunity to evaluator coupling ($\gamma{\approx}0.03$) on DeepSeek-chat notwithstanding, we observe that this immunity is model-dependent and can invert on newer model versions ($\gamma{\approx}1.06$ on DeepSeek-V4-Pro; descriptive---traces unarchived), highlighting the fragility of calibration-based mitigations.

\subsection{The Bias-Variance Tradeoff in Evaluation}

The coupling-noise tradeoff we identify echoes classical results in estimation theory. \citet{bubeck2012regret} formalize the exploration-exploitation tradeoff in online learning, where reducing bias (exploiting known-good strategies) increases variance (overfitting to evaluator noise). \citet{wolpert1997no}'s No Free Lunch theorem establishes that no algorithm can simultaneously optimize all performance metrics. Our contribution is to instantiate this abstract tradeoff in the concrete setting of LLM evaluator calibration, providing empirical measurements of the coupling-noise frontier and a formal theorem (Theorem~\ref{thm:tradeoff}) characterizing its boundary conditions.

\subsection{Measurement Reliability in LLM Evaluation}

The reliability paradox~\citep{messing2026tee}---that calibration can improve point estimates while worsening test-retest reliability---is closely related to our coupling-noise tradeoff. \citet{guo2017calibration} and \citet{tian2023just} study calibration in static evaluation settings, but do not consider closed-loop adaptation. Our work extends this literature to the TTRL setting, where evaluator calibration feeds back into agent strategy dynamics. The joint-measurement protocol we contribute (\S4.3) addresses the methodological challenge of simultaneously measuring coupling and noise without confounding.

\section{Discussion}

\subsection{Practical Guidelines}

\begin{enumerate}[leftmargin=*,nosep]
    \item \textbf{Limit communication rounds if temporal fatigue holds in your setting.} The reported decomposition makes temporal horizon the dominant cost, motivating $T \leq 10$ rounds; because we could not reproduce that decomposition, we present this as a precautionary default rather than an established guideline.
    \item \textbf{Don't trust self-evaluation.} Model updates can silently invert calibration behavior; validate periodically~\citep{zhou2024incontext}.
    \item \textbf{Use confidence-gating selectively.} It reduces coupling in the unarchived confidence-gating runs ($-37\%$; descriptive); the simulation grid does not show a robust noise increase (mean $\Delta\CV = -3.5\%$ under the deterministic seed); deploy when coupling is the binding constraint~\citep{chen2025unicr}.
    \item \textbf{Monitor the landscape.} No existing configuration occupies the favorable regime $\{\gamma < 0.2, \CV < 0.3\}$~\citep{caleffects2026}.
\end{enumerate}

\subsection{Calibration Monitor: A Practical Tool}\label{sec:monitor}

Based on our theoretical results (Theorems~\ref{thm:fatigue}--\ref{thm:tradeoff}), we release a \textbf{Calibration Monitor} tool that provides real-time diagnostics for multi-agent LLM systems. The tool computes three metrics per evaluation round:

\begin{table}[H]
\centering
\caption{\textbf{Calibration Monitor metrics and decision thresholds.}}
\label{tab:monitor}
\begin{tabular}{@{}lll@{}}
\toprule
Metric & Threshold & Action \\
\midrule
$\ECE > 0.20$ & Red & Reduce $T$ or restart agents \\
$\gamma > 0.50$ & Amber & Apply confidence-gating \\
$\CV > 0.30$ & Red & Increase evaluator diversity \\
$\text{CNR} < 2.0$ & Amber & Tradeoff unfavorable; skip calibration \\
\bottomrule
\end{tabular}
\end{table}

The decision tree is:
\begin{enumerate}[leftmargin=*,nosep]
    \item If $\ECE > 0.20$: calibration fatigue is dominant---reduce communication rounds.
    \item If $\gamma > 0.50$ and $\CV < 0.30$: coupling is the binding constraint---apply confidence-gating.
    \item If $\CV > 0.30$: noise is the binding constraint---increase evaluator diversity (TEE).
    \item If $\text{CNR} < 2.0$: the tradeoff is unfavorable---do not apply calibration.
\end{enumerate}

\noindent No anonymized source repository is provided: we do not claim a released implementation, and the checks above are fully specified by the metric definitions and thresholds in Table~\ref{tab:monitor}. Thresholds are heuristic defaults chosen to separate the operating regimes observed in \S3 (short- vs.\ long-horizon ECE bands; landscape $\gamma$/$\CV$ ranges) and require task-specific tuning (see Limitations).

\subsection{Cross-Domain Scope}

The calibration fatigue phenomenon (Finding~1) was reported in the math-reasoning domain (GSM8K) but is \emph{not} reproduced in our cross-model reruns. A code-generation (MBPP) pilot shows the same temporal-dominance pattern, but its per-run traces are not archived with this paper, so it is reported below as an unarchived pilot observation rather than a validated claim:

\begin{table}[H]
\centering
\caption{\textbf{Cross-domain calibration fatigue (GSM8K: reported, not replicated here; MBPP: pilot, unarchived).} The temporal-dominance pattern as originally reported; the MBPP row is a pilot observation. Neither row is reproduced in our cross-model reruns.}
\begin{tabular}{@{}lccc@{}}
\toprule
Domain & Task & $\Delta_{\text{pure } T}$ & $\Delta_{\text{pure sync}}$ \\
\midrule
Math reasoning & GSM8K & 90.8\% & 9.2\% \\
Code generation & MBPP & \multicolumn{2}{c}{\footnotesize pilot, unarchived} \\
\bottomrule
\end{tabular}
\end{table}

The MBPP pilot suggests a directionally larger communication effect, consistent with the hypothesis that code review-like interaction introduces additional calibration-relevant signals~\citep{sun2020test}; however, because neither the pilot traces nor precise estimates are archived, cross-domain generalizability of Theorem~\ref{thm:fatigue} remains future work~\citep{temporal2026}.

\subsection{Alternative Explanations}

We consider three alternative explanations that could account for the reported GPT-4o decomposition, and assess the evidence for each. Because that decomposition is itself unreplicated and its per-repetition traces are not archived, this analysis constrains how the reported claim should be interpreted rather than establishing its mechanism:

\begin{enumerate}[leftmargin=*,nosep]
    \item \textbf{Context accumulation.} Agents may accumulate their own prior outputs in the prompt, and this growing context---rather than time per se---drives degradation. Condition Q2 (Table~\ref{tab:results}) tests this by replicating C1b with agents blinded to their own prior rounds. Final ECE under no-history isolation is $0.208$, nearly identical to the with-history baseline C1b ($0.219$; 95\% CIs overlap), ruling out context accumulation as the primary driver.
    \item \textbf{Question memorization.} Fixed question pools may cause memorization that degrades calibration. We note that accuracy declines modestly ($0.69 \to 0.65$) while ECE increases disproportionately ($+0.21$), suggesting the effect is alignment-specific rather than a general capacity issue.
    \item \textbf{Evaluator version drift.} API updates may change evaluator behavior over time. We control for this by running all conditions within the same API version window, and the within-subjects design (each seed runs both standard and calibrated TTRL) eliminates inter-session drift.
\end{enumerate}

Pure temporal contrast is $\approx 0$ in both full-protocol replications (qwen2.5:32b $\delta=-0.004$, $p=0.86$; Qwen3.8-27B $\delta=+0.0003$, $p=0.99$), while the synchronization contrast is the sole significant driver on qwen2.5:32b ($+0.082$, $p=2\times10^{-5}$, $d=0.85$) but non-significant on Qwen3.8-27B ($+0.029$, $p=0.21$); the GPT-4o 90.8\% temporal attribution is thus model-specific, not model-general.

Calibration can drive ECE to the numerical floor ($\approx 10^{-9}$, isotonic, all $\alpha$/$\tau$) while maximizing coupling (conf\_gating $\gamma$ $0.29\to0.84$ as $\alpha$ $0\to2$ at $\tau=0.5$); $\tau$ dampens this ($0.84\to0.45$ from $\tau=0.5$ to $2.0$ at $\alpha=2.0$)---ECE-perfect does not imply coupling-safe.

All eight conditions sit $0.34$--$0.65$ above GPT-4o paper values (e.g., C1 $0.59$ vs $0.012$), indicating a baseline miscalibration offset rather than a condition-specific failure.

\subsection{Limitations}

(i) Finding~1's decomposition is reported from the original GPT-4o runs; three independent replications---deepseek-v4-flash (causal control) and full eight-condition, 30-seed qwen2.5:32b and Qwen3.8-27B reruns---did not reproduce the temporal fatigue effect (the latter with a reversed effect structure: synchronization significant, temporal null), so the finding is model-specific; this means our results may not generalize to models with fundamentally different calibration dynamics (e.g., open-source models with RLHF-trained evaluators). Per-rep traces for the GPT-4o Finding~1 runs are not included in the delivered evidence package; the headline statistics (including the $t$/$d$ statistics in \S3.1) should be read as reported-from-runs pending re-verification. We also record two reproducibility disclosures: the previously released CSV for this experiment was generated from the reported condition means rather than from raw run records, so it cannot serve as independent verification either; and no external timestamped registration exists, so we describe the $T{=}50$ strengthening criterion ($d > 0.30$) as \emph{pre-specified} in our working documents rather than preregistered. (ii) Finding~2 is based on $N{=}5$; the companion 30-seed expansion reports the same direction (31\% reduction; companion manuscript~\citep{companion2026}), but with $N{=}5$ the magnitude remains imprecise. Per-seed API traces for the Finding~2 runs are likewise not archived in the delivered package. \textbf{Ledger mapping:} per \texttt{CLAIM\_LEDGER.md}, C06 (90.8\% temporal), C07 ($\gamma{\approx}1.06$), C08 ($-42\%$ coupling), and C10 (35--37\% confidence-gated reduction) are gap/descriptive---reported from runs without archived per-repetition traces and treated as claims under test, not verified findings; only C09 (joint $2{\times}2$ null at $T{=}30$/$T{=}50$) and the simulation-grid means are verified from archived records. (iii) Finding~3's noise analysis is simulation-based~\citep{resampling2026}; real-API validation with diverse evaluator model pairs remains an open question. (iv) The confidence-gating update-rule confound (additive vs.\ multiplicative) is partially addressed by symmetric-LR controls but not fully isolated~\citep{boundarysync2026}. (v) Theorem~\ref{thm:tradeoff} assumes bounded directional signal strength; this may not hold for evaluators with very high position bias~\citep{tversky1974anchoring}. (vi) The calibration monitor thresholds are empirically derived and may require task-specific tuning~\citep{naeini2015obtaining,jiang2021can}. \textbf{Gap roadmap.} The falsification core is complete; what remains is all confirmatory and external-facing: a closed-source full-protocol replication with per-repetition archiving to close the verifiability gap, a two-family symmetric replication to test cross-family robustness of the null, a natural-cache ecological pilot to bound real-world leakage, and an extended-horizon joint run to adjudicate the delayed-communication hypothesis (planned protocols in \texttt{MERGE\_MAP\_A.md}~\S5; no new numbers are claimed here).

% ============================================================
% \S6 CONCLUSION
% ============================================================


\textbf{Reproducibility and independent verification.} Finding~1 (temporal calibration fatigue) is reported from the original GPT-4o runs (API version gpt-4o-2024-05-13, 30 repetitions per condition, deterministic seeds 1--30); those per-repetition traces are not archived, so its statistics cannot be recomputed (Limitations~(i)). An independent causal control on deepseek-v4-flash (three repetitions per condition) did \textit{not} reproduce the temporal fatigue effect: the C2 final ECE at $T{=}20$ overlapped the $T{=}10$ range, and the no-peer/placebo control ordering was inconsistent with the sync-driven mechanism (evidence: \texttt{p1\_causal\_quick\_20260808.json}, \texttt{p1\_calibration\_pressure\_probe\_20260808.json}). A full eight-condition, 30-paired-seed replication on qwen2.5:32b (240/240 runs, all per-round records released) again found no temporal effect ($-0.004$, $p = 0.86$) but a significant synchronization effect ($+0.082$, $d = 0.85$)---an effect-structure reversal relative to the GPT-4o decomposition; A second full eight-condition replication on Qwen3.8-27B (240/240 runs, per-round records released) again found a zero temporal contrast ($+0.0003$, $p = 0.99$) but \emph{no} significant synchronization effect, so that reversal is itself not model-general. All replication data ship with this paper and are recomputable.

\textbf{Reproduction failure analysis.} We investigate three potential causes for the reproduction failure: (i) \textit{Model architecture differences}: deepseek-v4-flash uses a hybrid attention mechanism with sliding-window attention for long contexts, which may reduce context-length calibration degradation compared to GPT-4o's full attention. (ii) \textit{Calibration priors}: deepseek-v4-flash may have been RLHF-trained with explicit calibration objectives, producing better-calibrated confidence estimates that are less susceptible to temporal drift. (iii) \textit{Sample size}: the causal control used only three repetitions per condition and is therefore underpowered on its own; with two full 30-paired-seed replications now showing a temporal contrast statistically indistinguishable from zero, however, the failure is not plausibly attributable to sample size. We conclude that Finding~1 is \textit{model-specific} and pending re-verification with GPT-4o-class access; cross-model generalization remains an open question. Our Theorem~\ref{thm:fatigue} provides a lower bound conditional on monotonic degradation, which may not hold for all model families.

\textbf{Sample size constraints.} Several conditions use $N{=}5$ or $N{=}8$ seeds, which limits the precision of coupling and CV estimates. The bootstrap CIs we report ($B{=}5000$) quantify sampling uncertainty but cannot address systematic bias from seed selection. The 30-seed companion experiment~\citep{companion2026} partially addresses this for the self-evaluation condition, but the full eight-condition design would benefit from $N{\geq}30$ per condition to establish robust effect sizes.

\textbf{Simulation grid limitations.} The coupling-noise simulation grid (Finding~3) uses a simplified TTRL dynamics model with fixed learning rates ($\alpha_{\text{win}}{=}0.08$, $\alpha_{\text{lose}}{=}0.04$) and a fixed strategy space ($n{=}11$). The mean $\Delta\CV{=}-3.5\%$ under the deterministic seed is seed-sensitive across script versions, so we cannot conclude that calibration universally reduces noise. Real API experiments with calibrated evaluators are needed to validate the simulation; our Theorem~\ref{thm:tradeoff} characterizes the conditions under which the tradeoff should hold, but empirical confirmation across diverse evaluator-model pairs remains pending.

\textbf{External validity.} All experiments use synthetic tasks (math, reasoning, classification) with artificial communication protocols. Whether temporal calibration fatigue generalizes to real-world multi-agent deployments---where communication is asynchronous, tasks are open-ended, and agents have memory beyond the current context window---is an open empirical question. The calibration monitor tool we contribute (\S\ref{sec:monitor}) provides a practical instrument for measuring fatigue in deployed systems, but longitudinal studies in production settings are needed to establish external validity.

\section{Conclusion}

The central result of this study is negative: the widely reported temporal calibration fatigue effect in multi-agent LLM systems does not replicate. Three independent replications across model families fail to reproduce it---the pure temporal contrast is statistically indistinguishable from zero in every one (point estimates $-0.004$ to $+0.0003$, $p \geq 0.86$)---and the original GPT-4o decomposition that motivated the claim is reported from runs whose per-repetition traces are not archived, so it cannot be recomputed from the package. We therefore report it as the claim under test rather than an established finding; the communicative component is not model-general either (it is significant in qwen2.5:32b but not in Qwen3.8-27B), so what generalizes is the absence of temporal fatigue rather than any positive driver.

What survives is narrower but fully archived. A deterministic simulation grid shows no universal noise increase from calibration (mean $\Delta\CV = -3.5\%$, seed-sensitive across script versions), so the coupling--noise tradeoff is better framed as a method-specific hypothesis---carrying explicit falsification conditions---than as a structural law; a companion 30-seed real-API expansion reports $\Delta\CV < 0$~\citep{companion2026}. Completed $2 \times 2$ joint-measurement experiments ($N{=}120$ seeds at $T{=}30$, extended by $N{=}120$ further seeds at $T{=}50$) show that neither Communication nor Calibration alters convergence outcomes: a suggestive late-phase Communication trend at $T{=}30$ fails seed-level testing and does not strengthen at $T{=}50$ (pre-specified criterion $d > 0.30$ unmet), leaving the delayed-communication hypothesis unconfirmed.

Taken together, these results define an empty favorable regime in the calibration landscape~\citep{cover2012elements}. Breaking the tradeoff---finding methods that reduce both coupling and noise simultaneously---remains an open problem, analogous to the bias-variance frontier in classical estimation theory~\citep{lehmann1998theory}. More broadly, our experience is a reminder that a striking cross-sectional decomposition can survive scrutiny while resting on data no reader can re-examine; we hope the fully archived replication core shipped here makes the negative result easier to build on than the original positive claim.

\section*{Broader Impact Statement and Ethical Considerations}

This work documents structural limitations of calibration in multi-agent large language model systems. Understanding these limitations is essential for safe and ethical deployment.

\textbf{Ethical implications of self-evaluation fragility.} Our Finding~2 demonstrates that model updates can silently invert self-evaluation immunity: a model version with near-zero coupling ($\gamma \approx 0.03$) can be replaced by an update with strong coupling ($\gamma \approx 1.06$), without any change in the evaluation protocol. This has direct ethical implications for autonomous systems that rely on self-evaluation for safety-critical decisions (e.g., medical triage, legal document review, autonomous driving). If a system's self-evaluation becomes coupled to its strategy distribution after a model update, it may become overconfident in flawed strategies while appearing calibrated under static testing. Practitioners must treat self-evaluation reliability as a version-dependent property that requires re-verification after every model update, not a one-time certification.

\textbf{Ethical implications of calibration fatigue.} Finding~1 reports that temporal calibration fatigue dominates calibration degradation in long-horizon multi-agent interactions (90.8\% of variance in the GPT-4o runs; model-specific, unverifiable from the published evidence, and \emph{not replicated in any of our three independent tests}). If that reported pattern holds in deployed systems with extended interaction horizons (e.g., customer service chatbots with multi-turn conversations, collaborative coding assistants), agents would become progressively miscalibrated without any external signal of degradation. Because the effect is simultaneously unreplicated and unverifiable from available data, we treat this as a \emph{risk hypothesis} rather than a demonstrated hazard---which is itself a reason for caution, since a widely cited hazard that no one can currently re-measure is hard to either confirm or dismiss. We therefore recommend that multi-agent LLM deployments include continuous calibration monitoring (see \S\ref{sec:monitor} for our contributed tool) and automatic circuit-breakers that trigger human review when ECE exceeds a threshold, independently of which mechanism turns out to drive degradation.

\textbf{Responsible research practices.} The central result of this paper is a reproduction failure, and we report it in full rather than suppressing it or selectively reporting only a model on which the original effect held. This transparency is essential for the research community to accurately assess the generalizability of the reported calibration-fatigue effect. All code, data, and analysis scripts for the replicable part of this work are released with deterministic seeds and SHA-256 checksums to enable independent verification; where evidence could not be archived, we mark the affected numbers as descriptive throughout rather than presenting them as verified. We do not make claims about real-world deployment safety beyond what our controlled experiments support; all external validity claims are explicitly flagged as open questions requiring further study.

\section*{Reproducibility Statement}

\textbf{Code and data availability.} The evidence package accompanying this paper is curated under \texttt{evidence/} and pinned by SHA-256 in \texttt{evidence/evidence\_manifest.json} (919 files). It contains: the joint $T{=}30$ and $T{=}50$ per-seed datasets (\texttt{f9\_uniform\_T30\_results/}, \texttt{f9\_t50\_results/}; 120 seeds each, 30 per condition) with their analysis summaries and generating scripts; the full eight-condition, 30-seed qwen2.5:32b and Qwen3.8-27B replications (\texttt{finding1\_rerun\_qwen32b/} and \texttt{finding1\_rerun\_qwen38\_27b/}, 240 runs each with per-round records); a small third-model pilot with its script (\texttt{finding1\_rerun\_qwen25\_7b\_pilot/}); the deterministic simulation grid (\texttt{canonical\_results\_v5.json}) with its script \texttt{evidence/canonical\_calibration\_v5.py}; the evaluator-landscape survey (\texttt{landscape\_5conditions.json}); and the Finding-1 data-gap documentation (\texttt{finding1\_data\_GAP.md}, \texttt{FINDING1\_DATA\_SEARCH\_VERDICT.md}). We do \emph{not} ship a consolidated \texttt{code/} or \texttt{data/} tree: the original GPT-4o run logs and scripts are not available to us, and the per-repetition artifact for that experiment was never archived (see below).

\textbf{Step-by-step reproduction protocol.} To reproduce the results this package does support: (1) \textit{Environment}: Python 3.10+ with \texttt{numpy} and \texttt{scipy>=1.10}; all offline analyses here are CPU-only. (2) \textit{Joint experiments}: the T30/T50 per-seed JSONs under \texttt{evidence/f9\_uniform\_T30\_results/} and \texttt{evidence/f9\_t50\_results/} are the single canonical source; \texttt{evidence/regen\_T30\_singlesource\_20260910.py} regenerates the T30 CSV and summary from them and reproduces the condition means and window tests in \S\ref{sec:joint}. (3) \textit{Simulation grid}: \texttt{evidence/canonical\_calibration\_v5.py} regenerates \texttt{canonical\_results\_v5.json} deterministically (per-method seed offsets; hash randomization disabled, \texttt{PYTHONHASHSEED=0}). (4) \textit{Cross-model replication}: the qwen2.5:32b contrasts are recomputable from the 240 per-round records in \texttt{evidence/finding1\_rerun\_qwen32b/}; \texttt{evidence/scripts/verify\_finding1.py} cross-checks them against the values reported here.

\textbf{Known reproduction limitations.} The GPT-4o decomposition that motivated this study \emph{cannot} be reproduced from this package: its per-repetition artifact was never archived, and we were unable to locate it across local disks, archives, and public repositories (attempt log: \texttt{evidence/FINDING1\_DATA\_SEARCH\_VERDICT.md}). For transparency we also record that the previously released CSV for that experiment was itself generated from the reported condition means rather than from raw run records, and therefore cannot serve as independent verification either. The deepseek-v4-flash causal control exists only as a small pilot (three repetitions per condition) and is reported descriptively throughout.

% ============================================================
% A-LINE APPENDIX STUBS (skeleton only; no new numbers).
% Appendix A: retrieval-side leakage (tmlr-p19 + f1-single cloud rerun).
% Appendix B: pilot explorations (tmlr-p9/p10 calibrated-TTRL + self-eval pilots; protocol deposit only).
% Appendix C: negative-result records (tmlr-p15 FROZEN declaration pointer; tmlr-p18 TARP failure).
% Full source mapping: CLAIM_LEDGER.md (per-file provenance; included in supplementary). Body kept frozen to avoid bloat.
% ============================================================

\bibliographystyle{tmlr}
\bibliography{references,refs_added}

\end{document}
